\chapter{De Rham cohomology}

\lecture{31}{Thu 15 Jan 2026 22:19}{Personal Review}

\section{The de Rham cohomology groups}

Recall that a closed form is a form $\omega $ such that $\mathrm{d} \omega =0$, and an exact form is any form of the form $\mathrm{d} \eta $. Since $\mathrm{d} ^2=0$, clearly exact implies closed. All closed forms are \emph{locally} exact, but not necessarily globally exact, meaning that the question of whether closed implies exact depends entirely on the global topological properties of the manifold.

Since $\mathrm{d} ^2=0$, defining $\Omega ^k(M)=0$ for all $k<0$ or $k>n$, then we have a cochain complex
\[\begin{tikzcd}[row sep = 2.5em, column sep = 2.5em]
\Omega ^0(M)\ar[" \mathrm{d}  ", r]  & \Omega ^1(M)\ar[" \mathrm{d}  ", r]  & \Omega ^2(M)\ar[" \mathrm{d}  ", r]  & \cdots
\end{tikzcd}\]
with values in $\mathbb{R} $-vector spaces. The $k$th cohomology group can then be described as the quotient of the set of closed $k$-forms by the set of exact $k$-forms, and thus closed implies exact if and only if the cochain complex has vanishsing homology in degree $k\geq 1$. The cohomology groups are called the \emph{de Rham cohomology groups}, and are denoted $H_{\mathrm{dR} } ^k(M)$.

\begin{lemma}\label{6.1.1}
	Let $F : M \to N$ be smooth. Since there exist maps $F^*: \Omega ^k(N) \to \Omega ^k(M)$ for all $N$, there exists a functor $\mathcal{M} \mathrm{an}^{\mathrm{op}} \to \Ch(\mathbb{R} ) $ given by mapping manifolds to their de Rham cochain complexes, and by mapping $F$ to the chain map given by pullbacks.
\end{lemma}
\begin{proof}
	Since pullbacks are functorial, it suffices to show naturality, meaning $F^*\circ \mathrm{d} =\mathrm{d} \circ F^*$. This follows by 3.1.22 (4).
\end{proof}
\begin{corollary}\label{6.1.2}
	The map $M\mapsto H_{\mathrm{dR} } ^k(M)$ is functorial $\mathcal{M} \mathrm{an} ^{\mathrm{op}} \to \mathbb{R} \text{-} \Vect$.
\end{corollary}
\begin{proof}
	Homology is functorial. In particular, any smooth $F : M \to N$ lowers to a map $F^*: H_{\mathrm{dR} } ^k(N) \to H_{\mathrm{dR} } ^k(M)$ defined by $F^*[\omega ]=[F^*\omega ]$.
\end{proof}
\begin{corollary}\label{6.1.3}
	Diffeomorphic manifolds have isomorphic de Rham cohomology groups. Even more powerfully, diffeomorphic manifolds have isomorphic (not just quasi-isomorphic) de Rham complexes of differential graded algebras.
\end{corollary}

\section{Elementary computations}

\begin{proposition}\label{6.2.1}
	Let $\left\{ M_j \right\} _{j\in J} $ be a countable collection of smooth manifolds. Then $H_{\mathrm{dR} } ^k$ preserves countable products in the sense that
	\[
		\prod_{j\in J} H_{\mathrm{dR} } ^k(M)\cong H_{\mathrm{dR} } ^k \coprod_{j\in J} M_j
	.\]
\end{proposition}
\begin{proof}
	It suffices to show $M\mapsto \Omega ^k$ preserves products in this sense, since limits in functor categories are computed pointwise, $\Ch(\mathbb{R} ) $ can be realized as a functor category, and homology preserves countable products. The isomorphism for $\Omega ^k$ is given by the inclusion maps $i_j : M_j\hookrightarrow \coprod_{j} M_j$ by
	\[
		\Omega ^k\coprod_{j\in J} M_j\ni\omega \mapsto \left\{ i^*_j\omega  \right\} _{j\in J} \in \prod_{j\in J} \Omega ^k(M_j)
	.\]
	I find this obvious and omit the proof.
\end{proof}

\begin{proposition}\label{6.2.2}
	Let $M$ be a connected smooth manifold. Then $H^0_{\mathrm{dR} } (M)$ is the space of constant functions, and thus 1-dimensional.
\end{proposition}
\begin{proof}
	There are no $(-1)$-forms, and thus there are no exact 0-forms. A 0-form is a smooth function $f : M \to \mathbb{R} $, and a closed 0-form is a function $f$ such that $\mathrm{d} f=0$. Recall this is true if and only if $f$ is constant.
\end{proof}

\begin{proposition}\label{6.2.3}
	Let $M$ be a 0-manifold. Then the de Rham cohomology of $M$ is concentrated in degree 0, and is a direct product of 1-dimensional vector spaces for each point of $M$.
\end{proposition}

\section{Homotopy invariance}

There is a strengthening of \cref{6.1.3}.
